\subsection{The $10 \times 10$ revival: a case study}
The $10 \times 10$ torus is the largest torus admitting the
longest period $N = 60$, and provides a clean illustration of
the revival mechanism. We take
$\rho_x = \rho_y = (5-\sqrt{5})/10 \approx 0.2764$ and
$\delta_x = \delta_y = 0$.

Figure~\ref{fig:10x10-fidelity} tracks the fidelity of the walker
against the initial state over eighty steps. The fidelity reaches
its maximum at $N = 60$ with a value numerically indistinguishable
from $1$; no other step in the range produces a full-state
revival. Figure~\ref{fig:10x10-probability} shows the position
probability distribution at ten snapshots. The walker starts
localised at the origin, spreads outward over the first thirty
steps, and refocuses to the origin at $N = 60$, matching the
initial condition almost exactly.

Figure~\ref{fig:10x10-spin} shows the four components of the spin
distribution over the same range. The coin state returns to its
initial values at both $N = 30$ and $N = 60$. The return at
$N = 30$ is a \emph{partial} revival: the coin reconstructs itself
while the position distribution has not yet refocused, and no
full-state revival occurs. Only at $N = 60$ do coin and position
both coincide with their initial values.
